\documentclass[10pt,a4paper]{article}
\usepackage[margin=24mm]{geometry}
\usepackage[T1]{fontenc}
\usepackage{lmodern}
\usepackage{amsmath,amssymb,booktabs,array}
\usepackage[hidelinks]{hyperref}
\setlength{\parskip}{4pt}
\setlength{\parindent}{0pt}
\emergencystretch=2em
\title{Adapting Native Geometry Decoders for Dynamic Actor Surfaces:\\
Interim Evidence on Coverage and Observed Free Space}
\author{WorldSim V7.3 research draft}
\date{Evidence snapshot: 8 September 2026}
\begin{document}
\maketitle
\begin{abstract}
We study metric surface reconstruction of rigid dynamic Actors from sparse
LiDAR and multi-camera observations, with known identities, calibration and
trajectories. The prototype adapts a pretrained native geometry decoder and
uses spatial queries to generate explicit local surfaces in Actor coordinates.
An interim development study shows that better coverage does not ensure
better physical first intersections: the completed joint candidate R5 increases
surface recall and reduces missing returns, but introduces more early surfaces
and observed-free-space intrusion than strong same-input references.
Changing the static background reduces a separate source of intrusion while
exposing a coverage cost. These observations motivate controlled comparisons
of label span, geometry adaptation and surface-based objectives. The report
records completed experiments and their limitations; matched native and joint
full-track training, the next objective comparison, and independent new-log
confirmation remain unfinished. No final method superiority is claimed.
\end{abstract}

\section{Task and research hypothesis}
For Actor $i$, we infer a reusable surface $\mathcal S_i$ in metric canonical
coordinates from a fixed build window. Known rigid motion places it at query
time: $\mathcal S_i(t)=T_{w\leftarrow i}(t)\mathcal S_i$. Camera calibration and
trajectories are read-only. We focus on rigid vehicle categories; learned
motion prediction, nonrigid deformation, intensity, complete no-return beam
simulation and photorealistic rendering are outside this experiment.

The hypothesis is that adapting a pretrained geometric decoding path and
jointly reading local multi-view and LiDAR evidence can improve surface
position and coverage without increasing violations of observed free space.
This is an empirical hypothesis, not a consequence of adding parameters or
attention. Earlier small-adapter results did not directly test it: their final
surface displacement was computed from a physical hidden state, while current
visual features affected an evidence update. In that interface, displacement
had zero direct derivative with respect to current visual features when
parameters and other inputs were fixed. V7.3 opens a trainable visual-to-surface
path; a negative outcome for one implementation cannot reject all geometric
foundation-model adaptation.

\section{Implemented model and supervision}
\subsection{Native decoding and spatial queries}
The official VGGT training framework already supports freezing the aggregator
while fine-tuning task heads~\cite{vggt}. Our prototype trains the full native
DPT depth head (32,654,562 parameters), reading four aggregation levels. Only
the completely frozen prefix is cached. No upper-aggregation LoRA is trained
in the present joint candidate. Direct native-depth supervision uses build
Actor LiDAR projected to the correct camera time, with the original metric
scale fixed by robust build-side alignment.

The query decoder has 1,670,517 parameters, hidden width 192 and three update
layers. Up to 1,024 build-evidence queries and 512 completion queries carry
canonical positions and learned states. Completion is seeded by aligned
native geometry, with a LiDAR fallback; the explicit image-only entry also
permits learned coarse positions when both point sources are absent. The
older main runs did not train that last input regime.
For a query $x_i$, view $v$ is read at
\begin{equation}
 x_{c,v}=T_{c\leftarrow w}(t_v)T_{w\leftarrow i}(t_v)x_i.
\end{equation}
Each query samples four local offsets at each of four feature scales, retaining
camera, time and direction information. This follows the reference-centered
sparse sampling idea of Deformable DETR~\cite{deformable}, transferred here
to calibrated Actor queries. Query chunks have size 128. Local
16-neighbor interactions use a spatial index with detached neighbor identities;
we do not construct a dense query--query distance matrix. This reduces the
interaction support to local neighborhoods; nearest-surface supervision still
costs $O(PF)$ for $P$ target points and $F$ faces, with chunked temporary arrays.
The prototype's image mask represents calibrated frustum and valid-image
support, not verified occlusion visibility.

Each query predicts an explicit $3\times3$ patch with eight triangles, normal
and bounded local bending, on a fixed 0.06\,m grid spacing. Query centers are
not limited to small displacements around original LiDAR candidates. Adjacent
vertices within a patch are connected, but different patches are not welded
into a closed manifold. Appearance opacity cannot hide a physical surface.

\subsection{One surface, distinct measurement semantics}
The implemented objectives act on this explicit geometry:
\begin{equation}
 L=L_{\mathrm{coverage}}+\lambda_n L_{\mathrm{native}}
       +\lambda_f L_{\mathrm{free}}+0.05L_{\mathrm{envelope}}
       +\lambda_e L_{\mathrm{event}}.
\end{equation}
Coverage is the mean distance from sampled measured target points to the nearest
triangle. Sparse targets do not define a complete surface: an unmeasured region
is not automatically empty, and we do not penalize all generated points merely
for lacking a nearby sparse target. Native supervision is axial-depth SmoothL1
with a 0.2\,m transition. It remains available at known build correspondences
even when current predicted geometry leaves the Actor support region. This
addresses a previously observed self-closing native-support path.

For the original beam with first measured range $d_r^*$ and predicted first
intersection $\hat d_r$, the hard free-space objective is
\begin{equation}
 L_{\mathrm{free}}=\frac1{|R|}\sum_{r\in R}
 \mathbf1[\hat d_r<\infty]\,[d_r^*-0.2-\hat d_r]_+.
\end{equation}
All near-box original beams contribute; an occluded rear Actor does not replace
the original first-return distance. Missing predictions have zero free penalty
and remain penalized through measured coverage and reported missing-return
rates. Regions behind the original return remain unknown. A finite beam-tube
version uses the same surface with fixed width 0.03\,m and $32\times32$ sampling;
it does not learn confidence or transparency.

The event term is a capped, fixed-footprint likelihood of the same first visible
surface with Gaussian range width 0.2\,m and cap 28. It is evaluated on the owned
subset of the same sampled beams. Repeating a surface does not add an independent
termination opportunity. Depth-termination supervision has prior art in
DS-NeRF~\cite{dsnerf}; changing an expected-depth loss to a likelihood is not
claimed as novel. Crucially, an entirely absent surface receives the cap and
zero position gradient from this term. Coverage and free-space geometry must
retain their separate roles. A local differentiable visibility boundary is not
a guarantee of global support discovery or convergence. The rasterized
visibility and antialiasing derivatives use nvdiffrast~\cite{nvdiffrast}.

\section{Data, comparators and estimands}
The nuScenes development study contains 489 rigid Actor entries in 31 windows
from 25 independent logs: 414 fit entries in 20 logs and 75 development entries
in five logs. The main runs admit 371 fit and 67 development entries with build
LiDAR; the remaining 43 and eight stay in the population with explicit missing
predictions. Each window supplies six cameras at four build times, jointly
encoded at $378\times672$. Two interleaved times supply evaluation returns.
Development logs never supply optimizer updates to the shared models.

Additional full-track labels are fit-only: 414 target files contain 1,690,284
measured Actor points, including positive targets for some empty-build entries.
Those labels never become prediction inputs. R5 and R6 use the short-window
label regime; R7--R9 use the longer fit-only labels. R10 and R11 are the matched
joint/native full-track comparison still running at this snapshot. R12 isolates
the beam-tube free objective relative to R10. An equal-capacity pointwise control
has not been trained on this full population; therefore a spatial-attention
causal claim is currently unsupported.

We retain LiDAR accumulation with PCA patches, native adapted geometry plus
canonical LiDAR fusion, native TSDF fusion, a CAPA-style build-only adaptation
and a trained AdaPoinTr reference. CAPA supplies an official sparse-measurement
test-time adaptation formulation~\cite{capa}; our transfer resets its low-rank
parameters per window, takes 100 build-only updates and fuses its outputs in
canonical coordinates. Its per-image affine alignment and test-time budget
differ from shared training. AdaPoinTr~\cite{adapointr} uses the official PCN
initialization, 30 fit epochs, axis conversion and full-track labels. It has a
different optimizer, pretraining corpus and output density. These are meaningful
competitors, not all single-factor ablations.

An owned original return is a hit within $\pm0.2$\,m, early below that interval,
late above it, or missing if no surface intersects. These outcomes include
misses in the denominator. Free-space intrusion uses all original near-box
beams. Measured surface recall is the fraction of heldout Actor measurements
within 0.2\,m of the explicit surface. Frames are weighted by their actual
measurements within each Actor, followed by Actor means within logs and equal
log weighting. Paired uncertainty resamples independent logs 10,000 times,
not individual rays. With only five development logs, intervals remain
limited evidence. Twenty-three development Actors have no owned heldout
return; they are retained and their unavailable metrics are identified.

\section{Completed development evidence}
\begin{table}[htbp]
\centering\small
\begin{tabular}{lrrrrr}
\toprule
Method & Hit (\%) & Early (\%) & Miss (\%) & Free (m) & Recall (\%) \\
\midrule
LiDAR PCA & 21.24 & 4.35 & 73.41 & 0.0177 & 65.42 \\
Actor TSDF + carving & 3.28 & 0.26 & 96.33 & 0.0008 & 21.51 \\
Native fusion & 22.78 & 9.16 & 57.35 & 0.1503 & 78.13 \\
LiDAR R6 (window labels) & 31.01 & 14.02 & 49.31 & 0.0838 & 72.42 \\
LiDAR R7 (track labels) & 31.75 & 11.84 & 49.90 & 0.0716 & 74.55 \\
LiDAR R8 (beam free) & 31.00 & 5.59 & 56.83 & 0.0343 & 71.95 \\
LiDAR R9 (+event) & 32.47 & 6.01 & 54.57 & 0.0385 & 72.47 \\
CAPA-style TTA R2 & 17.80 & 10.95 & 58.46 & 0.1203 & 76.90 \\
AdaPoinTr R2 & 14.29 & 8.54 & 54.24 & 0.1221 & 84.60 \\
Joint R5 (window labels) & 22.28 & 20.98 & 37.71 & 0.3517 & 79.71 \\
\bottomrule
\end{tabular}

\caption{Completed candidates, all 75 development entries in five logs.
Recall is measured-target-to-surface recall, not complete-surface recall.
Empty predictions remain missing. Short/long labels, adaptation budgets and
output representations differ as described in the text; no row is a final
confirmed method. Late rates are the residual outcome and are omitted for space.}
\label{tab:population}
\end{table}

R5 completes 30 effective epochs (11,130 updates) of genuine native DPT and
query adaptation. The final recall is 79.71\%, but early returns reach 20.98\%
and mean free intrusion reaches 0.3517\,m. Relative to its own initialization,
missing returns decrease while early returns increase. Relative to LiDAR R6,
the surface is closer to heldout measurements but physical first intersections
are worse (Table~\ref{tab:paired}). These are not contradictory metrics:
additional support can cover true measurements and still lie in front of other
observations. The main candidate has not established a physical advantage.

\begin{table}[htbp]
\centering\small
\begin{tabular}{llrr}
\toprule
Joint R5 minus & Metric & Difference & 95\% log bootstrap \\
\midrule
R5 initialization & Early (pp) & +8.706 & [+1.221, +18.493] \\
R5 initialization & Miss (pp) & -13.433 & [-16.301, -8.843] \\
LiDAR R6 & Hit (pp) & -8.728 & [-14.097, -4.478] \\
LiDAR R6 & Free (m) & +0.268 & [+0.056, +0.503] \\
Native fusion & Early (pp) & +11.826 & [+5.400, +23.016] \\
Native fusion & Free (m) & +0.201 & [+0.037, +0.400] \\
\bottomrule
\end{tabular}

\caption{Selected paired development differences, five independent logs.
Positive early/free and negative hit are adverse. R5 versus R6 shares the
short-window label regime, but differs in visual input and native adaptation;
it does not isolate attention. The recovered R5 run also has the random-state
limitation described in Section~\ref{sec:limits}.}
\label{tab:paired}
\end{table}

A separate observed-return point evaluation uses only original owned beam
intersections and measured targets, with missing predictions retained. R5's
F-score is 0.3541, compared with R6's 0.4319 and R9's 0.4922. R5 minus R6 is
$-0.0778$ with log-bootstrap interval $[-0.1173,-0.0382]$. This is an observable
subset, not full canonical-surface precision. Unsquared two-direction Chamfer
is conditional on nonempty predicted/measured sets and is never substituted
for the missing-rate column. Nearest-neighbor matching can also hide a
wrong-beam intersection; literal first-return metrics remain necessary.

The R7--R9 sequence isolates narrower mechanisms on LiDAR queries. Longer fit
labels reduce early/free relative to R6. The beam-tube free objective then
reduces early/free further while lowering measured recall. Adding event
supervision reduces missing returns relative to R8, but the hit/free paired
intervals span zero; event support is still absent on roughly 39\% of sampled
owned training beams. None of these results proves that event likelihood alone
recovers missing geometry.

Query-source accounting attributes 91.74\% of R5's log-weighted free intrusion
and 78.37\% of its early-return contribution to completion-initialized patches.
Their centers average 1.087\,m from the nearest original build point. This
measures source provenance, not causal visual attribution or displacement
from each query's initial position: both source groups pass through the same
trainable spatial/visual updates.

\section{Density, static background and composition}
Point-cloud conversion changes the physical readout. AdaPoinTr R2 retains all
16,384 generated points, but its main comparison uses a common limited set of
patch centers. Reading PCA patches at all native output points increases
development hit by 24.72 percentage points and early by 30.85 points, while
increasing free intrusion by 0.2190\,m. A denser surface is not an unqualified
improvement. Native Actor TSDF at 0.1\,m voxels and 0.3\,m truncation produces
a nonempty mesh for only 162 of 489 entries (18 of 75 development entries).
Its low free intrusion therefore coexists with 96.33\% missing returns.
This is an applicability limit of the stated sparse configuration, not a
general negative result about TSDF fusion.

At scene level we preserve background/Actor ownership and cast against
\begin{equation}
 \mathcal S_{\mathrm{scene}}(t)=\mathcal S_{\mathrm{static}}
      \cup\bigcup_i T_{w\leftarrow i}(t)\mathcal S_i.
\end{equation}
The closest actual intersection across all owners is the scene first return.
All methods share background inputs, Actor trajectories and the readout.
Static construction excludes all known boxes at each build observation time.
Original build beams provide one-pass face carving before the measured return;
this does not fill unobserved background behind Actors.

Using VDBFusion~\cite{vdbfusion} with fixed 0.1\,m voxels, 0.3\,m truncation and
no hole filling yields a substantially different background support. On the
five-log nuScenes development background-only comparison, TSDF plus carving
reduces mean free intrusion from 0.2656 to 0.0822\,m, but increases missing
returns from 53.28\% to 62.80\%. Every log improves free intrusion and loses
hit coverage. With this common TSDF background, R5 still exceeds R9 in Actor
cohort early returns by 15.33 points and free intrusion by 0.1084\,m; both
paired intervals exclude zero. Background correction does not remove the
Actor geometry problem.

One already-exposed AV2 development log provides a separate implementation
check. Its points are already compensated into reference-ego coordinates;
endpoints are transformed to world once, while each return retains its actual
time-dependent sensor origin. The single-origin VDBFusion API is called in
groups of exactly equal origins, with no time/origin averaging. Relative to
carved PCA background, carved TSDF reduces background-only free intrusion from
0.5130 to 0.1075\,m but increases missing from 32.13\% to 45.54\%. Changing only
the background leaves R5's moving-Actor outcomes unchanged on its 88 moving
cohort beams. This single development log cannot supply between-log uncertainty
or independent confirmation.

\section{Remaining experiments and claim boundaries}\label{sec:limits}
R10 and R11 still require completed training and analysis. R12 is registered
but unrun. The one-epoch R13 input diagnostic is registered for all 51 original
zero-build-LiDAR entries, with 41 camera-supported fit entries trained from
M1 initialization. Seven have no positive full-track target but retain original
near-box rays; unknown space must not become a negative target. A single epoch
will establish execution and gradient behavior, not sufficient fitting or
generalization. Its expanded input condition is separate from the objective
comparison. The eight development empty-build entries have only one owned
heldout return, preventing a credible geometric accuracy conclusion for that
subgroup.

Twenty additional AV2 logs were selected by available build metadata before
model quality evaluation. Their 936 Actor entries, raw beams and 28-view
frozen prefixes are prepared; TSDF background construction is in progress.
No model-quality confirmation has been run on these logs. The final method
and background choice must precede that readout. This is cross-dataset
confirmation at a different camera/resolution budget, not an i.i.d. nuScenes
test. Project-unseen logs also do not prove absence from foundation-model
pretraining.

Current joint and native-only configurations run on the RTX 3090 while
preserving the stated time/view inputs. The joint peak allocated memory is approximately
10.2\,GiB; native-only training uses approximately 2.35\,GiB. Twenty-eight-view
AV2 prefix extraction at $672^2$ peaks at 7.64\,GiB. These are operation-specific
allocated peaks, not a sum or a bound for deeper PEFT. Container RAM is capped
at 90\,GiB. Concurrent run wall time is not exclusive GPU time.

The original R5 process disappeared after epoch 21 without an available
traceback; the cause remains unknown. Recovery restored model and optimizer
but its legacy checkpoint lacked RNG state. The recovered run is not a
bitwise continuation: 107 incomplete-epoch updates were discarded, and the
observed combined wall-time lower bound is 9.32 hours. The run provenance
remains part of its interpretation.

Annotation boxes with margins approximate point ownership and overlap regions
are excluded; they are not dense instance ground truth. nuScenes here lacks
per-return timestamps, whereas the AV2 adapter has per-return sensor origins
with a documented, development-inferred dual-LiDAR ID convention. Frustum
support is not occlusion truth. The learned patches are not globally manifold.
Trajectory-edit examples demonstrate composition only: with no counterfactual
ground truth, released unknown background is not a correctness measurement.
These limitations, the small development-log count and the unfinished controls
preclude a final reconstruction or sensor-simulation superiority claim.

\section*{Evidence and reproducibility}
The source branch is \texttt{research/worldsim-v7.3-geometric-adaptation},
derived directly from the final V7.2 state. The current operational authority
is \path{docs/RESEARCH_STATUS.md}; \path{docs/RESEARCH_FAILURES.md} and
\path{docs/EXPERIMENTS.md} retain failures and experiment history. Tables in
this draft are generated from saved JSON summaries, without rerunning
inference or selecting observations. Raw run directories under
\path{/root/autodl-tmp/runs/worldsim_v73/} retain configurations, checkpoints,
training logs and explicit surfaces. \path{paper_v73/README.md} identifies
the exact compact evidence files and compilation commands. This draft is
not an arXiv submission and does not mark V7.3 as complete.

\begin{thebibliography}{9}
\bibitem{vggt} VGGT authors. Official training framework and native geometry
implementation. \url{https://github.com/facebookresearch/vggt/tree/main/training}.
\bibitem{adapointr} X. Yu et al. AdaPoinTr: Diverse Point Cloud Completion with
Adaptive Geometry-Aware Transformers. IEEE TPAMI, 2023.
\url{https://arxiv.org/abs/2301.04545}.
\bibitem{capa} CAPA authors. Depth Completion as Parameter-Efficient Test-Time
Adaptations. 2026. Official project and implementation:
\url{https://research.nvidia.com/labs/dvl/projects/capa/}.
\bibitem{dsnerf} K. Deng et al. Depth-Supervised NeRF: Fewer Views and Faster
Training for Free. CVPR, 2022. \url{https://www.cs.cmu.edu/~dsnerf/}.
\bibitem{vdbfusion} I. Vizzo et al. VDBFusion: Flexible and Efficient TSDF
Integration of Range Sensor Data. Sensors, 2022.
\url{https://github.com/PRBonn/vdbfusion}.
\bibitem{deformable} X. Zhu et al. Deformable DETR: Deformable Transformers
for End-to-End Object Detection. 2020. Official implementation:
\url{https://github.com/fundamentalvision/Deformable-DETR}.
\bibitem{nvdiffrast} S. Laine et al. Modular Primitives for High-Performance
Differentiable Rendering. ACM TOG 39(6), 2020.
\url{https://nvlabs.github.io/nvdiffrast/}.
\end{thebibliography}
\end{document}
